\mychapter{Partie II : Réseaux Convolutifs (CNN)}{Partie II : Réseaux Convolutifs (CNN)}

\section{Contexte et jeu de données}
\slidepurpose{Exploiter l'indépendance spatiale des données visuelles avec les CNN.}{Partie 2 / Section 1}

La deuxième partie s'intéresse à la vision par ordinateur avec le jeu de données CIFAR-10. Contrairement aux données tabulaires de la première partie, les images possèdent une forte corrélation spatiale locale. CIFAR-10 contient 60 000 images couleur (RGB) de dimensions 32x32 réparties en 10 classes (avions, voitures, chats, etc.). Les données ont été prétraitées en tenseurs PyTorch et normalisées pour chaque canal.

\section{Compréhension des opérations fondamentales}

Afin de maîtriser les mécanismes internes des CNN, nous avons implémenté "from scratch" les opérations fondamentales de corrélation croisée 2D et de pooling (Max et Average) avant de les valider formellement contre l'API native de PyTorch.

\textbf{Implémentation manuelle de la corrélation 2D (\texttt{corr2d}) :}
\begin{verbatim}
def corr2d(X, K):
    h, w = K.shape
    Y = torch.zeros((X.shape[0] - h + 1, X.shape[1] - w + 1))
    for i in range(Y.shape[0]):
        for j in range(Y.shape[1]):
            Y[i, j] = (X[i:i+h, j:j+w] * K).sum()
    return Y
\end{verbatim}

Une comparaison directe avec \texttt{nn.Conv2d}, \texttt{nn.MaxPool2d} et \texttt{nn.AvgPool2d} via \texttt{torch.allclose} a permis de valider mathématiquement la conformité de nos boucles locales.

\section{Architecture LeNet5CIFAR}

Nous avons implémenté une architecture de référence inspirée du célèbre LeNet-5, modifiée pour admettre des images RGB et intégrant des améliorations modernes telles que la \texttt{BatchNorm2d} et le \texttt{Dropout}.

\begin{itemize}
    \item \textbf{Bloc 1 :} Conv2d (3$\rightarrow$16 filtres, 5x5, padding=2) $\rightarrow$ BatchNorm $\rightarrow$ ReLU $\rightarrow$ MaxPool2d(2x2)
    \item \textbf{Bloc 2 :} Conv2d (16$\rightarrow$32 filtres, 5x5, padding=0) $\rightarrow$ BatchNorm $\rightarrow$ ReLU $\rightarrow$ MaxPool2d(2x2)
    \item \textbf{Classifieur :} Flatten $\rightarrow$ LazyLinear(120) $\rightarrow$ ReLU $\rightarrow$ Dropout(0.3) $\rightarrow$ Linear(84) $\rightarrow$ Dropout(0.2) $\rightarrow$ Linear(10)
\end{itemize}

Après un entraînement sur 10 époques (optimiseur Adam, taux d'apprentissage de 0.001), la précision finale (\textbf{Accuracy}) sur les 10 000 images de test a atteint \textbf{72.93\,\%}. La perte finale s'est établie à 0.7681.

\solutionfigure{figures/fig_p2_train_curves.png}{Évolution de la Perte et de la Précision}

\section{Étude d'ablation avec \texttt{FlexibleCNN}}

Afin de quantifier l'impact de chaque choix de design, un modèle entièrement paramétrable (\texttt{FlexibleCNN}) a été conçu. Six scénarios ont été entraînés de façon systématique sur 3 époques pour comparer leur convergence initiale.

\begin{center}
\renewcommand{\arraystretch}{1.3}
\begin{tabular}{l c}
\rowcolor{TableHeader} \tblhead{Configuration Architecturale} & \tblhead{Précision Test (\%)} \\
\rowcolor{TableAlt} \textbf{Baseline (Standard)} ($p=2, s=1$, MaxPooling, 16 filtres) & 67.97\,\% \\
\textbf{Sans Padding} ($p=0$) & 67.05\,\% \\
\rowcolor{TableAlt} \textbf{Stride Fort} ($s=2$) & 63.54\,\% \\
\textbf{Average Pooling} & 66.86\,\% \\
\rowcolor{TableAlt} \textbf{Plus de Filtres} (64 filtres) & 71.21\,\% \\
\textbf{Avec Conv 1x1} (Compression des canaux) & 70.45\,\% \\
\end{tabular}
\end{center}

\begin{alertblock}{Analyse Architecturale}
Les résultats soulignent l'importance du MaxPooling pour préserver les gradients forts comparativement à l'Average Pooling. L'augmentation du nombre de filtres (capacité du réseau) apporte un net gain (71.21\%), tout comme l'introduction de convolutions 1x1 (Bottleneck structurel, 70.45\%), démontrant l'utilité des mélanges de canaux intra-pixels. Un stride fort en début de réseau provoque en revanche une perte d'information préjudiciable (baisse à 63.54\%).
\end{alertblock}

\section{Visualisation des Cartes de Caractéristiques (Feature Maps)}

Pour comprendre l'extraction hiérarchique, le modèle en mode évaluation a été appliqué à une image de la classe \emph{cat} (chat).

\solutionfigure{figures/fig_p2_cat_original.png}{Image originale dé-normalisée}

Les activations en sortie de la première couche convolutionnelle (16 filtres) et de la seconde couche (32 filtres) révèlent la transition des caractéristiques géométriques bas niveau (contours) vers des motifs plus abstraits.

\solutionfigure{figures/fig_p2_feature_maps_conv1.png}{Feature maps de la première couche de convolution (Conv1)}
\solutionfigure{figures/fig_p2_feature_maps_conv2.png}{Feature maps de la deuxième couche de convolution (Conv2)}
